% ============================================================
\setpartpreamble{\partintro
  {Operating in production: Row Level Security, tenant-safe relationships, retrofitting,
  performance, and the way from the plugin into the application and into CI.}
  {
    \item Chapter~\ref{ch:row-level-security-in-production}: Row Level Security in
    production
    \item Chapter~\ref{ch:tenant-safe-relationships}: Tenant-safe relationships
    \item Chapter~\ref{ch:retrofitting-and-operating}: Retrofitting and operating
    multi-tenancy
    \item Chapter~\ref{ch:plugin-performance}: Plugin performance
    \item Chapter~\ref{ch:using-the-plugin-in-asp}: Using the plugin in ASP.NET Core
    \item Chapter~\ref{ch:ci-and-api-stability}: CI and API stability
  }}
\part{Operation and integration}\label{part:operations}

% ------------------------------------------------------------
\chapter{Row Level Security in production}\label{ch:row-level-security-in-production}
\versionrisk{stable}

RLS is only as strong as its operating conditions. These pitfalls silently render the
policies ineffective or unusable:

\begin{datatable}{H{0.93}H{0.93}H{1.14}}
  \textbf{Pitfall} & \textbf{Consequence} & \textbf{Remedy} \\
  \tablehead
  App runs as table owner & Policies are ignored & \code{FORCE ROW LEVEL SECURITY} (plugin
  default) and separate roles \\
  \tablerule
  Superuser or \code{BYPASSRLS} & RLS is always bypassed, even with FORCE & App role with
  \code{NOSUPERUSER NOBYPASSRLS} (custom operation, chapter~\ref{ch:custom-migration-operations}) \\
  \tablerule
  Setting not set & Without \code{nullif}: error; with it: empty results & The interceptor
  sets the value on every open \\
  \tablerule
  No index on the tenant column & Every query scans the whole table & Index convention
  (chapter~\ref{ch:conventions}) \\
  \tablerule
  Transaction pooling (PgBouncer) & Session settings reach other clients or get lost &
  \code{Use\allowbreak{}Transaction\allowbreak{}Pooling(\allowbreak{})}: the setting
  transaction-locally (see below) \\
  \tablerule
  Connection opened outside EF & The session keeps an old tenant or none & Let EF open
  connections, or \code{Apply\allowbreak{}Tenant\allowbreak{}Session(\allowbreak{})}
  (chapter~\ref{ch:interceptors}) \\
  \tablerule
  “Bypass” switch for the app role & Anyone who can run SQL as the app role can flip it &
  Separate roles for cross-tenant access (see below) \\
\end{datatable}

\section{The threat model}\label{sec:threat-model}

RLS is a layer of protection with clear limits. Whoever uses it should know what it protects
against, what it relies on and where it ends:

\begin{datatable}{H{0.45}H{1.55}}
  \textbf{Aspect} & \textbf{In the sample} \\
  \tablehead
  Protected & The rows of tenant-owned tables in the supported mappings
  (page~\pageref{sec:supported-table-mappings}). Not protected are shared master data, the
  schema and, without composite keys, the existence of foreign values
  (chapter~\ref{ch:tenant-safe-relationships}). \\
  \tablerule
  Trusted & The database server and its roles; the migration role as table owner; the code of
  the application and the plugin, including the interceptors; the
  \code{ITenant\allowbreak{}Provider} and the source it reads the tenant from. \\
  \tablerule
  Assumptions & The app role is neither superuser nor owner and has no \code{BYPASSRLS}. The
  tenant comes from an authenticated claim, not from a header or parameter the client chooses
  freely. EF opens the connections, and there is no transaction pooling, except with
  \code{Use\allowbreak{}Transaction\allowbreak{}Pooling(\allowbreak{})}. \\
  \tablerule
  Limit & The app role may set the session setting itself; it has to, so that the interceptor
  can set it. Whoever runs arbitrary SQL under this role switches the tenant with it. \\
\end{datatable}

The limit is not a flaw of the plugin but a property of the session model.
\code{Threat\allowbreak{}Model\allowbreak{}Tests} shows it against PostgreSQL: one
\code{SELECT set\_\allowbreak{}config(\allowbreak{}'app.\allowbreak{}tenant\_\allowbreak{}id',
\allowbreak{} .\allowbreak{}.\allowbreak{}.\allowbreak{})} under the app role is enough, and the
query after it returns the other tenant's rows.
\code{Sql\allowbreak{}Server\allowbreak{}Threat\allowbreak{}Model\allowbreak{}Tests} shows the
same with \code{sp\_\allowbreak{}set\_\allowbreak{}session\_\allowbreak{}context}. This defines
what RLS protects against here and what it does not:

\begin{itemize}
  \item \textbf{Protected:} forgotten filters, faulty LINQ queries, the application's raw SQL
    with a wrong \code{WHERE}, bulk operations, reports that accidentally run through the app
    role. In all these cases the plugin sets the tenant, and the code does not change it.
  \item \textbf{Not protected:} SQL injection and compromised application code. Both can change
    the setting, even within a single statement, since \code{set\_\allowbreak{}config} is a
    function. What helps against them are parameters instead of concatenated strings (EF
    parameterizes LINQ and \code{From\allowbreak{}Sql}/\code{Sql\allowbreak{}Query} with
    interpolated strings; \code{From\allowbreak{}Sql\allowbreak{}Raw} with concatenated values
    does not) and privileges that give the app role only the tables it needs.
  \item \textbf{Outside the database:} a forged or wrongly assigned claim grants access to
    exactly that tenant. That is the job of authentication; the plugin takes the tenant from
    the \code{ITenant\allowbreak{}Provider}.
\end{itemize}

The sample deliberately does not implement possible hardenings. SQL Server can freeze the key
until the end of the connection with \code{sp\_\allowbreak{}set\_\allowbreak{}session\_\allowbreak{}context
.\allowbreak{}.\allowbreak{}.\allowbreak{}, @read\_only = 1}
(\href{https://learn.microsoft.com/en-us/sql/relational-databases/security/row-level-security}{SQL Server: Row-Level Security});
the test above shows that a later change then fails. The price: \code{Apply\allowbreak{}Tenant\allowbreak{}Session},
too, can no longer change the tenant of an open connection. PostgreSQL has no read-only setting;
there remain a role per tenant with \code{SET ROLE} or a signed tenant token checked by a
\code{SECURITY DEFINER} function. Both move complexity into the database and are only named as a
direction here. Cross-tenant access still belongs in roles of its own (section below).

\section{Transaction pooling}\label{sec:transaction-pooling}

PgBouncer in “transaction” mode shares a server connection between clients per transaction.
A session value set with
\code{set\_\allowbreak{}config(\allowbreak{}.\allowbreak{}.\allowbreak{}.\allowbreak{},\allowbreak{}
false)} can then reach another client, a security hole. The safe variant sets the value
\emph{transaction-locally} with
\code{set\_\allowbreak{}config(\allowbreak{}.\allowbreak{}.\allowbreak{}.\allowbreak{},\allowbreak{}
true)}: it ends with commit or rollback. For that, the interceptor moves from
\code{IDb\allowbreak{}Connection\allowbreak{}Interceptor} to
\code{IDb\allowbreak{}Transaction\allowbreak{}Interceptor} and sets the tenant in
\code{Transaction\allowbreak{}Started}, in a command of the new transaction. The sample offers
this as an option:

\begin{sourcecode}[Application]{csharp}
services.AddDbContext<ShopContext>((sp, options) => options
    .UseNpgsql(connectionString)
    .UseAcmeTenancy(sp, t => t.UseTransactionPooling()));
\end{sourcecode}

% source: src/Acme.EntityFrameworkCore.Tenancy.PostgreSQL/Runtime/PostgreSqlTransactionSessionInterceptor.cs
\begin{sourcecode}[Runtime/\allowbreak{}Postgre\allowbreak{}Sql\allowbreak{}Transaction\allowbreak{}Session\allowbreak{}Interceptor.\allowbreak{}cs (\allowbreak{}shortened)]{csharp}
public override async ValueTask<DbTransaction> TransactionStartedAsync(
    DbConnection connection,
    TransactionEndEventData eventData,
    DbTransaction result,
    CancellationToken cancellationToken = default)
{
    await using var command = CreateSessionCommand(connection, result, eventData.Context);
    await command.ExecuteNonQueryAsync(cancellationToken);
    return result;
}

// SaveChanges always in a transaction, even for a single command
private static void RequireTransaction(DbContextEventData eventData)
{
    // ...
        context.Database.AutoTransactionBehavior = AutoTransactionBehavior.Always;
}

protected virtual DbCommand CreateSessionCommand(
    DbConnection connection, DbTransaction transaction, DbContext? context)
{
    // ...
    var command = connection.CreateCommand();
    command.Transaction = transaction;
    command.CommandText = "SELECT set_config(@name, @value, true)";   // is_local = true
    // ...
}
\end{sourcecode}

\code{Use\allowbreak{}Acme\allowbreak{}Tenancy(\allowbreak{}sp)} then registers this interceptor
instead of the connection interceptor. Three rules come with it:

\begin{itemize}
  \item \textbf{Reading in transactions.} Queries outside a transaction have no tenant and see
    no rows. That is safe, but the application has to put reading work that needs data into a
    transaction.
  \item \textbf{Saving always in a transaction.} EF omits the automatic transaction when
    \code{Save\allowbreak{}Changes} runs only one command
    (\code{Auto\allowbreak{}Transaction\allowbreak{}Behavior.\allowbreak{}When\allowbreak{}Needed}).
    A single insert would then run without a tenant and fail on the policy. The interceptor is
    therefore also an \code{ISave\allowbreak{}Changes\allowbreak{}Interceptor} and sets
    \code{Always}.
  \item \textbf{No session value.} If the option is set and a
    \code{Tenant\allowbreak{}Connection\allowbreak{}Interceptor} is registered anyway, it refuses
    to set the value with a \code{Tenancy\allowbreak{}Exception} instead of leaving a value that
    outlives the transaction. That also rules out
    \code{Apply\allowbreak{}Tenant\allowbreak{}Session}: a call fails with a message that the
    tenant is set per transaction and the commands belong in one.
\end{itemize}

\code{Transaction\allowbreak{}Pooling\allowbreak{}Tests} checks against PostgreSQL: inside the
transaction the value is set and RLS applies, after commit and after rollback it is empty on the
same connection, a single insert gets its transaction, the connection interceptor refuses
the session value, and \code{Apply\allowbreak{}Tenant\allowbreak{}Session} explains why it is not available. PgBouncer itself is not part of the test; what is checked is the property
that matters: that no value remains after the transaction.

\section{Policy performance}

PostgreSQL inserts policies into every query like additional \code{WHERE} conditions.
\code{current\_\allowbreak{}setting} is \code{STABLE}. That guarantees the same result for
all rows within one statement and \emph{allows} the optimizer to combine several calls into
one (\href{https://www.postgresql.org/docs/17/xfunc-volatility.html}{PostgreSQL: Function Volatility Categories}); it does not guarantee a fixed
number of calls. How often the function actually runs depends on the plan. If the comparison
is part of an \code{Index Cond}, PostgreSQL evaluates the comparison value once and reads only
the tenant's rows. If it is a \code{Filter} over a seq scan, it may be evaluated per row, and
then the scan weighs more than the function. Expensive policies come from subqueries (for
example permissions from another table). In addition, PostgreSQL holds back conditions of the
query whose functions are not \code{LEAKPROOF} behind the policy conditions, which can prevent
optimizations (\href{https://www.postgresql.org/docs/17/ddl-rowsecurity.html}{PostgreSQL: Row Security Policies}).

Only the plan under real conditions is therefore meaningful: \code{EXPLAIN
(\allowbreak{}ANALYZE,\allowbreak{} BUFFERS)} as the app role with a tenant set, with a
realistic amount of data and current statistics. Measured as owner or superuser, the policies
are missing from the plan, because superusers and roles with \code{BYPASSRLS} bypass them, and
so does the owner without \code{FORCE}. The overhead shows in comparing the same query with
and without the policy, not in the number of function calls.

\section{Legitimate cross-tenant access}

Sooner or later someone needs data across tenants: reports for the operator, maintenance
jobs, support. The tempting shortcut is a switch in the application, for example an
additional session setting \code{app.\allowbreak{}bypass} that the policy honors. That is
not a security boundary: whatever the application role can set, an injected SQL statement
can set too. Separation must come from database roles instead:

\begin{datatable}{H{0.75}H{1.25}H{1.00}}
  \textbf{Need} & \textbf{Solution} & \textbf{Why it is safe} \\
  \tablehead
  Reports, analytics & Separate read-only role with \code{BYPASSRLS}, only \code{SELECT},
  its own connection string and a context configuration without tenant & The app role cannot
  become this role; its credentials live elsewhere \\
  \tablerule
  Support reads selected data & Additional permissive policy for a support role:
  \code{Has\allowbreak{}Policy(\allowbreak{}.\allowbreak{}.\allowbreak{}.\allowbreak{})}
  with \code{For(\allowbreak{}Select)} and
  \code{To(\allowbreak{}\textquotedbl{}shop\_\allowbreak{}support\textquotedbl{})} &
  Restricted to SELECT and to members of the role \\
  \tablerule
  Maintenance per tenant & Background job that iterates the tenants and sets each one
  (chapter~\ref{ch:using-the-plugin-in-asp}) & Runs through the normal policies \\
\end{datatable}

\begin{sourcecode}[Reporting role (\allowbreak{}tests/\allowbreak{}.\allowbreak{}.\allowbreak{}.\allowbreak{}Integration\allowbreak{}Tests/\allowbreak{}Cross\allowbreak{}Tenant\allowbreak{}Access\allowbreak{}Tests.\allowbreak{}cs,\allowbreak{} simplified)]{sql}
-- Created by an administrator: BYPASSRLS requires a superuser
CREATE ROLE shop_reporting LOGIN PASSWORD '...' NOSUPERUSER BYPASSRLS;
GRANT USAGE ON SCHEMA public TO shop_reporting;
GRANT SELECT ON ALL TABLES IN SCHEMA public TO shop_reporting;
\end{sourcecode}

The application role, in contrast, cannot switch the policies off. \code{SET
row\_\allowbreak{}security = off} does not disable RLS for a role without \code{BYPASSRLS};
PostgreSQL instead rejects every query the policies would affect.
\code{Cross\allowbreak{}Tenant\allowbreak{}Access\allowbreak{}Tests} checks all three: the
reporting role sees all tenants, it cannot write, and the app role cannot help itself.

\begin{important}[Audit cross-tenant access]
  Every connection of the reporting role is an access to all tenants. Log it
  (\code{log\_\allowbreak{}connections}, pgaudit) and grant the role to as few people and
  services as possible. The plugin's \code{Create\allowbreak{}Application\allowbreak{}Role}
  deliberately always creates roles with \code{NOBYPASSRLS}: bypass roles do not belong in
  the application's migrations.
\end{important}

% ------------------------------------------------------------
\chapter{Tenant-safe relationships}\label{ch:tenant-safe-relationships}
\versionrisk{stable}

Which table mappings the plugin protects and which it rejects while the model is built is set by
the matrix on page~\pageref{sec:supported-table-mappings}; this chapter does not repeat it. It is
about the second axis: relationships between protected tables. RLS governs which \emph{rows} a
session sees. Whether a row only refers to rows of the same tenant is not governed by RLS. That
is \emph{referential tenant consistency}, and it comes from the model.

The sample uses customer, order and item: \code{Customer} → \code{Order} →
\code{Order\allowbreak{}Item}, all three \code{ITenant\allowbreak{}Owned}, plus a shared
\code{Product} without a tenant.

\section{Visibility is not consistency}

PostgreSQL documents explicitly that checks of foreign keys, unique and primary keys bypass RLS
to maintain data integrity, and warns of “covert channels” through these checks
(\href{https://www.postgresql.org/docs/17/ddl-rowsecurity.html}{PostgreSQL: Row Security Policies}).
The SQL Server documentation says nothing explicit about it;
\code{Sql\allowbreak{}Server\allowbreak{}Threat\allowbreak{}Model\allowbreak{}Tests} shows that
the foreign key check there also sees rows the filter predicate hides. With foreign keys on the
ID alone, as EF creates them by default, this follows:

\begin{datatable}{H{0.70}H{1.30}}
  \textbf{Observation} & \textbf{Cause} \\
  \tablehead
  Tenant A adds an item to an order of B & The FK check finds the order, although A cannot
  see it. \\
  \tablerule
  A learns whether a foreign ID exists & Unknown IDs fail on the foreign key, foreign ones
  do not. \\
  \tablerule
  The own item vanishes from \code{Include(\allowbreak{}i => i.\allowbreak{}Order)} & A required
  navigation becomes an \code{INNER JOIN}, and the order is invisible. \\
  \tablerule
  A deletes an own order, an item of B disappears & The cascading delete is a referential action
  and bypasses RLS as well. \\
  \tablerule
  A learns that an e-mail address exists at B & A global unique index reports the
  collision. \\
\end{datatable}

\code{Relationship\allowbreak{}Tests} proves every row against PostgreSQL with a restricted
role. None of these cases is a read leak in the strict sense; A sees no foreign row. But the
data loses its tenant boundary, and the error messages answer questions A should not be able
to ask.

\section{Composite foreign keys}

The solution is to make the tenant part of every reference. The principal gets an alternate key
of tenant and ID, the dependent side a foreign key of its tenant column and the ID:

% source: tests/Acme.EntityFrameworkCore.Tenancy.IntegrationTests/RelationshipTests.cs
\begin{sourcecode}[tests/\allowbreak{}.\allowbreak{}.\allowbreak{}.\allowbreak{}Integration\allowbreak{}Tests/\allowbreak{}Relationship\allowbreak{}Tests.\allowbreak{}cs › Tenant\allowbreak{}Safe\allowbreak{}Context]{csharp}
modelBuilder.Entity<Customer>(b =>
{
    b.HasAlternateKey(c => new { c.TenantId, c.Id });
    b.HasIndex(c => new { c.TenantId, c.Email }).IsUnique();   // unique per tenant
});

modelBuilder.Entity<Order>(b =>
{
    b.HasAlternateKey(o => new { o.TenantId, o.Id });
    b.HasOne(o => o.Customer).WithMany(c => c.Orders)
        .HasForeignKey(o => new { o.TenantId, o.CustomerId })
        .HasPrincipalKey(c => new { c.TenantId, c.Id })
        .OnDelete(DeleteBehavior.Restrict);                   // customers are not cascaded
});

modelBuilder.Entity<OrderItem>()
    .HasOne(i => i.Order).WithMany(o => o.Items)
    .HasForeignKey(i => new { i.TenantId, i.OrderId })
    .HasPrincipalKey(o => new { o.TenantId, o.Id })
    .OnDelete(DeleteBehavior.Cascade);                         // items belong to the order
\end{sourcecode}

Now an item of A refers to the pair (A, ID). An order of B never satisfies it, and the error is
the same as for an unknown ID: there is nothing left to learn. The tenant column is part of the
foreign key and the tenant property at the same time. The SaveChanges interceptor sets it on new
items, and fixup through the navigation sets the same value. No index is needed for the policy,
because key and foreign key start with the tenant column; the index convention recognizes that.

The price is an additional unique index per principal and a wider foreign key. The plugin does
not create these keys automatically. A convention could, but it would have to rewrite every
relationship between protected types and reconcile it with explicit configuration. The sample
leaves the decision to the application's model.

\section{Uniqueness per tenant}

A unique index over a business column applies to the whole table, that is across all tenants.
Uniqueness per tenant is almost always what is meant: \code{(TenantId, Email)} instead of
\code{Email}. Where uniqueness really has to be global, for example for the tenants'
subdomains, the column belongs in the tenant directory, which deliberately holds no tenant data.

\section{Shared master data}

\code{Product} has no tenant, no policy and is readable by everyone. Foreign keys from tenant data
to master data are unproblematic; the test reads the product through the item. The opposite
direction, master data with references to tenant data, must not exist. Writing to master data
should not be allowed to the app role; \code{Grant\allowbreak{}Table\allowbreak{}Privileges}
grants privileges on all tables of a schema, a finer grant is up to the application's
migrations.

\section{Join tables}

The plugin rejects implicit many-to-many tables when one side is tenant-owned. A join entity with
a tenant property of its own is protected (page~\pageref{sec:supported-table-mappings}). For
consistency it gets the same composite foreign keys as every dependent table:
\code{(TenantId, LeftId)} and \code{(TenantId, RightId)} on the alternate keys of both sides.
Then a link can only connect two rows of the same tenant.

\section{Delete rules}

Cascading deletes run in the database as referential actions and bypass RLS. With composite
foreign keys they still stay within the tenant, because every dependent row carries the same
tenant. Without them a cascade deletes foreign rows; the test
\code{Naive\allowbreak{}\_cascade\allowbreak{}\_deletes\allowbreak{}\_rows\allowbreak{}\_of\allowbreak{}\_another\allowbreak{}\_tenant}
shows exactly that. For roots such as the customer, \code{Restrict} is usually right, for parts of
an aggregate such as items \code{Cascade}. EF also cascades in the change tracker, but only for
loaded entities; the rule in the database is what counts.


% ------------------------------------------------------------
\chapter{Retrofitting and operating multi-tenancy}\label{ch:retrofitting-and-operating}
\versionrisk{stable}

The previous chapters assume a new database. Usually, however, there already are tables with
data, a running application and users who keep working during the change. This chapter shows the
way from a table without a tenant column to RLS, and after that operation under retries and lost
connections. All steps are executable in \code{Operations\allowbreak{}Tests} against PostgreSQL.
The chapter closes with deployment: migrations in the pipeline with separate roles, and changes
without downtime.

\section{From a table without a tenant to RLS}

Whoever simply lets an existing entity implement \code{ITenant\allowbreak{}Owned} gets a single
migration: a column with the empty tenant as default, RLS, policy, index. After it nobody sees the
old rows any more. A backfill in the same migration does not reliably help either: with
\code{FORCE ROW LEVEL SECURITY} the policy also applies to the migration role, and without a tenant
set its \code{UPDATE} finds no row. The way therefore leads through two migrations, following the
pattern expand, fill, protect:

\begin{datatable}{H{0.22}H{0.90}H{1.88}}
  \textbf{No.} & \textbf{Step} & \textbf{Approach} \\
  \tablehead
  1 & Expand & The entity implements \code{ITenant\allowbreak{}Owned} but stays excluded with
  \code{Has\allowbreak{}No\allowbreak{}Row\allowbreak{}Level\allowbreak{}Security(\allowbreak{})}.
  The migration only adds the column. \\
  \tablerule
  2 & Fill & As owner, in the same or a migration of its own: every row gets its tenant from
  whatever knows the assignment, for example a mapping table or the tenant of the parent row. \\
  \tablerule
  3 & Switch the application & An app version whose model no longer has the exclusion: the
  interceptor sets the tenant on new rows. The database does not protect yet; the model does not
  check that. \\
  \tablerule
  4 & Close the gap & Assign rows that were created without a tenant between steps 1 and 3. \\
  \tablerule
  5 & Protect & Second migration: exclusion removed, so \code{ENABLE}/\code{FORCE}, policy and
  index. First of all a guard that stops as long as a row has no tenant. \\
  \tablerule
  6 & Tighten & Optionally composite keys and unique indexes per tenant
  (chapter~\ref{ch:tenant-safe-relationships}). \\
\end{datatable}

The guard is a step of its own at the start of \code{Up}, before the operations EF generates:

\begin{sourcecode}[Guard (\allowbreak{}SQL)]{postgresql}
DO $$ BEGIN
  IF EXISTS (SELECT 1 FROM "Tickets"
             WHERE "TenantId" = '00000000-0000-0000-0000-000000000000') THEN
    RAISE EXCEPTION 'Tickets without a tenant: backfill first';
  END IF;
END $$;
\end{sourcecode}

\begin{sourcecode}[Migrations/\allowbreak{}…\allowbreak{}\_Enable\allowbreak{}Ticket\allowbreak{}Isolation.\allowbreak{}cs (example)]{csharp}
protected override void Up(MigrationBuilder migrationBuilder)
{
    // Before anything else: do not protect without a complete assignment
    migrationBuilder.Sql(TicketGuardSql);      // the guard above, as a constant

    // ... generated by EF: ENABLE/FORCE ROW LEVEL SECURITY, CREATE POLICY, CREATE INDEX
}
\end{sourcecode}

The migration runs in a transaction. If the guard fires, the table stays without RLS, and no
half-done state arises;
\code{Guard\allowbreak{}\_stops\allowbreak{}\_enabling\allowbreak{}\_rls\allowbreak{}\_while\allowbreak{}\_rows\allowbreak{}\_lack\allowbreak{}\_a\allowbreak{}\_tenant}
checks this. After a complete assignment every tenant sees exactly its rows.

\subsection*{App versions, Down and rollback}

During the change two app versions often run at the same time, and each meets one of the two
states of the database:

\begin{datatable}{H{0.95}H{1.00}H{1.05}}
  \textbf{Database} & \textbf{App without tenant} & \textbf{App with tenant} \\
  \tablehead
  After migration 1 (column, no RLS) & Runs; new rows without a tenant (step 4) & Runs, sets
  the tenant; not yet protected \\
  \tablerule
  After migration 2 (RLS) & Fail closed: reads nothing, writes fail & Protected \\
\end{datatable}

The old app version must therefore no longer run once migration 2 is applied. Rolling the
application back to a version without a tenant thus first requires the Down of migration 2. The
Down removes policy and RLS; column and values stay, but the table is unprotected again. The test
checks both. A Down of migration 1 drops the column and with it the assignment; it belongs in
development only.

Policy changes follow the same pattern. If the setting name changes
(\code{Use\allowbreak{}Setting\allowbreak{}Name}), the migration replaces the policy in its
transaction. Between the migration and the new app version the policy reads a setting the running
app does not set, and the database refuses everything until both match. A transition in which the
interceptor sets both settings is conceivable but not implemented. For SQL Server,
page~\pageref{sec:sqlserver-session-settings} describes the same case.

\section{Retries, transactions and lost connections}

The plugin sets the tenant when EF opens a connection. This determines the behavior under EF's
resilience mechanisms:

\begin{itemize}
  \item \textbf{Retries.} An execution strategy such as
    \code{Enable\allowbreak{}Retry\allowbreak{}On\allowbreak{}Failure} repeats an operation after
    a transient error on a newly opened connection. The interceptor runs again, the tenant is set.
    The test ends the session on the server with \code{pg\_\allowbreak{}terminate\_\allowbreak{}backend};
    the next query fails, is repeated and returns only rows of the tenant.
  \item \textbf{Explicit transactions.} With a retrying strategy EF refuses commands in a
    transaction the application began itself, because a retry would have to repeat the whole
    transaction. The transaction therefore belongs into
    \code{Create\allowbreak{}Execution\allowbreak{}Strategy(\allowbreak{}).\allowbreak{}Execute\allowbreak{}Async(\allowbreak{}…)}.
    Inside, EF opens the connection once, sets the tenant and keeps it until the commit. If the
    connection breaks in the middle of the transaction, the strategy repeats the whole delegate on
    a new connection; side effects outside the database then have to be idempotent.
  \item \textbf{Session reinitialization.} Npgsql resets a connection's state when it returns to
    the pool, as long as \code{No Reset On Close} is not enabled
    (\href{https://www.npgsql.org/doc/connection-string-parameters.html}{Npgsql: Connection String Parameters}).
    SqlClient also resets the session context when it reuses a pooled connection
    (\code{sp\_\allowbreak{}reset\_\allowbreak{}connection}). The plugin relies
    on neither, but sets the tenant on every open; only connections that EF does not open need
    \code{Apply\allowbreak{}Tenant\allowbreak{}Session} (chapter~\ref{ch:interceptors}).
\end{itemize}

For transaction pooling, \code{Use\allowbreak{}Transaction\allowbreak{}Pooling(\allowbreak{})}
sets the tenant per transaction (page~\pageref{sec:transaction-pooling}). Not implemented is a
read-only session key (page~\pageref{sec:threat-model}). \code{Operations\allowbreak{}Tests} contains the lost connection and
the transaction inside the strategy; the cases of open connections are in
\code{Connection\allowbreak{}Lifetime\allowbreak{}Tests}.


\section{Deployment: pipeline, roles and zero downtime}\label{sec:deployment}

The samples create the database in development with \code{Development\allowbreak{}Database}.
In production, every schema change belongs in a pipeline that strictly separates two roles:

\begin{itemize}
  \item \textbf{\code{shop\_\allowbreak{}owner}} owns schema, tables and policies and runs the
    migrations.
  \item \textbf{\code{shop\_\allowbreak{}app}} is the application's role: only
    \code{SELECT}, \code{INSERT}, \code{UPDATE}, \code{DELETE}, always under the policies.
\end{itemize}

The application therefore never calls
\code{db.\allowbreak{}Database.\allowbreak{}Migrate(\allowbreak{})} on startup. Its role would
need DDL rights for that, and whoever may alter a table may also run
\code{DISABLE ROW LEVEL SECURITY} or drop the policy. \code{FORCE} binds the owner to the policy
but does not protect the policy from the owner. A hole in the application, such as an SQL
injection, would then no longer be contained by the database. Operations add to that: in a
rolling deployment several instances start at once, and a failed migration should stop the
rollout, not keep the instances in a restart loop.

\subsection*{Once: the roles}

Before the first pipeline runs, the database administrators create the roles. That is not a
migration, because migrations already run as the owner. The script is in the repository:

% source: docs/deployment/setup-roles.sql
\begin{sourcecode}[docs/\allowbreak{}deployment/\allowbreak{}setup-\allowbreak{}roles.\allowbreak{}sql]{postgresql}
CREATE ROLE shop_owner LOGIN NOSUPERUSER NOBYPASSRLS PASSWORD :'owner_password';
CREATE ROLE shop_app LOGIN NOSUPERUSER NOBYPASSRLS PASSWORD :'app_password';
CREATE DATABASE shop OWNER shop_owner;

...
GRANT USAGE ON SCHEMA public TO shop_app;

ALTER DEFAULT PRIVILEGES FOR ROLE shop_owner IN SCHEMA public
    GRANT SELECT, INSERT, UPDATE, DELETE ON TABLES TO shop_app;
ALTER DEFAULT PRIVILEGES FOR ROLE shop_owner IN SCHEMA public
    GRANT USAGE, SELECT ON SEQUENCES TO shop_app;
\end{sourcecode}

The elision is \code{\textbackslash{}connect shop}: from there on, the script runs in the new
database. The \code{DEFAULT PRIVILEGES} matter: every table that a later migration creates as
\code{shop\_\allowbreak{}owner} is usable by \code{shop\_\allowbreak{}app} right away. Without
them, the new application version fails with \code{permission denied}, although migration and
policy are correct. The same holds for sequences: an \code{INSERT} into a table with a
\code{serial} column calls \code{nextval} on its sequence, and EF's HiLo pattern does so before
the \code{INSERT}; without \code{USAGE} on the sequence the app role gets \code{permission
denied for sequence}. An identity column (\code{GENERATED ... AS IDENTITY}, Npgsql's default)
needs no privilege of its own on its internal sequence for an \code{INSERT}. The plugin's
\code{Grant\allowbreak{}Table\allowbreak{}Privileges} operation also sets default privileges
for tables and sequences, but only for objects that the role running it creates later; the
script sets them for \code{shop\_\allowbreak{}owner} independently of who runs it.

\subsection*{The pipeline}

\begin{flowdiagram}
  \flowstep{Build and test}{One commit, one build: tests, bundle and application are produced
  together}
  \flowstep{Migration bundle}{\code{dotnet ef migrations bundle}, self-contained; plus the
  idempotent script for review}
  \flowstep{Approval}{GitHub environment \code{production}: approval and secrets only for the
  following jobs}
  \flowstep{Bundle as shop\_owner}{Applies pending migrations in transactions; the only job
  with the owner connection}
  \flowstep{Rolling deployment}{Only after a successful migration; the application connects as
  \code{shop\_\allowbreak{}app}}
  \flowstep{Contract (later)}{A follow-up migration cleans up when no old instance is running
  any more}
\end{flowdiagram}

The template \code{docs/\allowbreak{}deployment/\allowbreak{}migrations-\allowbreak{}deploy.\allowbreak{}yml}
implements this with three jobs. It is not active in the repository; an application copies it
to \code{.github/\allowbreak{}workflows} and replaces the deploy step with its platform's:

% source: docs/deployment/migrations-deploy.yml
\begin{sourcecode}[docs/\allowbreak{}deployment/\allowbreak{}migrations-\allowbreak{}deploy.\allowbreak{}yml (\allowbreak{}shortened)]{text}
jobs:
  build:
    steps:
      ...
      - name: Migration bundle
        run: >
          dotnet ef migrations bundle --project "$APP_PROJECT" --context "$DB_CONTEXT"
          --configuration Release --self-contained -r linux-x64 -o artifacts/efbundle
      ...
      - run: dotnet publish "$APP_PROJECT" -c Release -o artifacts/app

  migrate:
    needs: build
    environment: production
    steps:
      ...
      - name: Apply pending migrations as shop_owner
        env:
          OWNER_CONNECTION: ${{ secrets.SHOP_OWNER_CONNECTION_STRING }}
        run: |
          chmod +x artifacts/efbundle            # artifacts lose the executable bit
          artifacts/efbundle --connection "$OWNER_CONNECTION"

  deploy:
    needs: migrate
    environment: production
    ...
\end{sourcecode}

The bundle contains the startup project with the plugin; for its options, what section
“Migration bundles” in chapter~\ref{ch:advanced-migrations} says applies.
\code{--connection} overrides the design-time factory's connection, and a second run applies
nothing. Only the \code{migrate} job knows the owner connection string; the application's is
given to the hosting platform, not to the pipeline.

The consumer test (\code{eng/\allowbreak{}consumer-\allowbreak{}test/\allowbreak{}run.\allowbreak{}sh
--nativeaot}, page~\pageref{sec:consumer-test}) runs exactly this sequence against PostgreSQL in Docker: it runs
\code{setup-\allowbreak{}roles.\allowbreak{}sql} unchanged, builds the bundle from the generated
migration, applies it twice as \code{shop\_\allowbreak{}owner} (the second run changes
nothing), checks that \code{shop\_\allowbreak{}app} may not alter the table
(\code{must be owner}), and starts the application as \code{shop\_\allowbreak{}app}, which sees
only its tenant's rows.

\subsection*{Zero downtime under RLS: expand and contract}

In a rolling or blue-green deployment, the migration runs before the new version, and for a
while old and new instances work against the same schema. Every migration must therefore work
with the previous version of the application. The expand-and-contract pattern splits a change
into three phases for that, like the conversion in section~“From a table without a tenant to
RLS”:

\begin{datatable}{H{0.70}H{1.15}H{1.15}}
  \textbf{Phase} & \textbf{Schema} & \textbf{What to watch under RLS} \\
  \tablehead
  1 Expand & New columns nullable or with a default, new tables, new indexes; nothing is
  removed or renamed. & New tables with \code{ITenant\allowbreak{}Owned} get their policy right
  away; the old version does not use them. If a policy has to change, a second one is added
  that allows the new form. \\
  \tablerule
  2 Roll out & Unchanged. The new version writes to the new columns, old instances drain. &
  Both versions set the same tenant; the policies must admit the rows of both versions. \\
  \tablerule
  3 Contract & Follow-up migration: drop old columns, set \code{NOT NULL}, remove temporary
  policies. & Only when no old instance is running; before that a guard as in section~“From a
  table without a tenant to RLS”. \\
\end{datatable}

For policies, a PostgreSQL property applies: several \code{PERMISSIVE} policies of a table are
combined with \code{OR}. An additional policy thus widens what is visible and serves as the
transition in phase~1; removing it in phase~3 is the contraction. Each of these policies must
separate tenants on its own, otherwise the transition opens the table. Because DDL is
transactional in PostgreSQL, replacing a policy within one migration leaves no moment without a
policy.

Rolling back the application is possible without a database change until phase~3, because the
schema still supports the old version. After that it needs the Down of the follow-up migration,
and dropped columns do not come back with it. A Down therefore does not belong in the pipeline;
an error after phase~3 is fixed forward with a new migration.

% ------------------------------------------------------------
\chapter{Plugin performance}\label{ch:plugin-performance}
\versionrisk{stable}

Plugin code runs at three frequencies. Knowing which code runs how often tells you where
optimization pays off:

\begin{datatable}{H{0.75}H{1.25}H{1.00}}
  \textbf{Frequency} & \textbf{Code} & \textbf{Consequence} \\
  \tablehead
  Once per internal provider or model & Conventions, type mapping plugins, singleton options & Readability
  over micro-optimization \\
  \tablerule
  Once per query shape & Translators, query interceptors & Moderate; the result is cached
  (chapter~\ref{ch:the-query-pipeline}) \\
  \tablerule
  On every execution & Interceptors for commands, connections, SaveChanges, materialization;
  evaluatable filter when the funcletizer extracts the parameter values & Hot path: no reflection, no unnecessary allocations \\
\end{datatable}

This leads to a few rules the sample follows: \code{Method\allowbreak{}Info} and
\code{Property\allowbreak{}Info} lookups live in static fields. The SaveChanges interceptor
only collects entries when there are tenant-relevant changes. The connection interceptor
costs one round trip per open; with many small queries in one request the application
therefore opens the connection once with
\code{Open\allowbreak{}Connection\allowbreak{}Async(\allowbreak{})}, and “one context, one
tenant” applies (chapter~\ref{ch:interceptors}). Such assumptions should be measured, not
guessed. The repository contains a BenchmarkDotNet project; the benchmark for the
SaveChanges logic runs without a database:

% source: benchmarks/Acme.EntityFrameworkCore.Tenancy.Benchmarks/SaveChangesInterceptorBenchmark.cs
\begin{sourcecode}[benchmarks/\allowbreak{}.\allowbreak{}.\allowbreak{}.\allowbreak{}/\allowbreak{}Save\allowbreak{}Changes\allowbreak{}Interceptor\allowbreak{}Benchmark.\allowbreak{}cs (\allowbreak{}shortened)]{csharp}
[MemoryDiagnoser]
public class SaveChangesInterceptorBenchmark
{
    [Params(1, 100)]
    public int Entities { get; set; }

    // GlobalSetup: context with UseAcmeTenancy() (no connection) and the interceptor

    [Benchmark(Baseline = true)]
    public int TrackOnly()                  // what EF does anyway before saving
    {
        AddOrders();
        var count = _db.ChangeTracker.Entries().Count();
        _db.ChangeTracker.Clear();
        return count;
    }

    [Benchmark]
    public int TrackAndApplyTenant()        // the same plus the plugin's SaveChanges logic
    {
        AddOrders();
        _interceptor.ApplyTenant(_db);
        var count = _db.ChangeTracker.Entries().Count();
        _db.ChangeTracker.Clear();
        return count;
    }
}
\end{sourcecode}

\begin{sourcecode}[Running the benchmark]{text}
dotnet run -c Release --project benchmarks/Acme.EntityFrameworkCore.Tenancy.Benchmarks \
    -- --filter '*SaveChanges*' --job short
\end{sourcecode}

\begin{datatable}{H{1.67}H{0.83}H{0.83}H{0.83}H{0.83}}
  \textbf{Method} & \textbf{Entities} & \textbf{Mean} & \textbf{Ratio} & \textbf{Allocated} \\
  \tablehead
  \code{Track\allowbreak{}Only} & 1 & 1.5 µs & 1.00 & 1.8 KB \\
  \tablerule
  \code{Track\allowbreak{}And\allowbreak{}Apply\allowbreak{}Tenant} & 1 & 2.3 µs & 1.54 &
  2.7 KB \\
  \tablerule
  \code{Track\allowbreak{}Only} & 100 & 123 µs & 1.00 & 134 KB \\
  \tablerule
  \code{Track\allowbreak{}And\allowbreak{}Apply\allowbreak{}Tenant} & 100 & 170 µs & 1.38 &
  202 KB \\
\end{datatable}

Example results (short job, .NET 10, author's machine; absolute values vary, the ratios are
more meaningful). The plugin adds about half a microsecond per saved entity, roughly 40 to
50 percent on top of change tracking alone. That sounds like much, but a single round trip
to the database costs hundreds of microseconds; the connection interceptor's extra round
trip per open (chapter~\ref{ch:interceptors}) therefore weighs far more than the SaveChanges
logic. That is the kind of insight a benchmark provides: it shows where optimizing is
worthwhile and where it is not.

\section{Reading and query compilation}

Two more benchmarks cover the other frequencies.
\code{Materialization\allowbreak{}Interceptor\allowbreak{}Benchmark} measures what
\code{Verify\allowbreak{}Materialized\allowbreak{}Tenants(\allowbreak{})} costs when reading.
It needs real rows from a \code{Db\allowbreak{}Data\allowbreak{}Reader} and starts PostgreSQL
with Testcontainers (Docker) for that. The connection stays open, so the connection
interceptor's round trip is not measured; \code{As\allowbreak{}No\allowbreak{}Tracking} keeps
change tracking out. A third variant with an empty
\code{IMaterialization\allowbreak{}Interceptor} separates EF's cost from the plugin's:

% source: benchmarks/Acme.EntityFrameworkCore.Tenancy.Benchmarks/MaterializationInterceptorBenchmark.cs
\begin{sourcecode}[benchmarks/\allowbreak{}.\allowbreak{}.\allowbreak{}.\allowbreak{}/\allowbreak{}Materialization\allowbreak{}Interceptor\allowbreak{}Benchmark.\allowbreak{}cs (shortened)]{csharp}
[MemoryDiagnoser]
public class MaterializationInterceptorBenchmark
{
    [Params(10, 100, 1000)]
    public int Rows { get; set; }

    // GlobalSetup: container, 1000 rows of one tenant, connections open

    [Benchmark(Baseline = true)]
    public Task<int> Read() => CountAsync(_plain);

    [Benchmark]
    public Task<int> ReadIntercepted() => CountAsync(_intercepted);

    [Benchmark]
    public Task<int> ReadAndVerify() => CountAsync(_verifying);

    private async Task<int> CountAsync(BenchmarkContext db)
        => (await db.Orders.AsNoTracking().OrderBy(o => o.Id).Take(Rows).ToListAsync()).Count;

    // Create: UseAcmeTenancy(sp, t => t.VerifyMaterializedTenants(verify))
}
\end{sourcecode}

\begin{datatable}{H{1.67}H{0.83}H{0.83}H{0.83}H{0.83}}
  \textbf{Method} & \textbf{Rows} & \textbf{Mean} & \textbf{Ratio} & \textbf{Allocated} \\
  \tablehead
  \code{Read} & 10 & 200 µs & 1.00 & 12 KB \\
  \tablerule
  \code{Read\allowbreak{}Intercepted} & 10 & 210 µs & 1.08 & 21 KB \\
  \tablerule
  \code{Read\allowbreak{}And\allowbreak{}Verify} & 10 & 220 µs & 1.12 & 22 KB \\
  \tablerule
  \code{Read} & 100 & 278 µs & 1.00 & 36 KB \\
  \tablerule
  \code{Read\allowbreak{}Intercepted} & 100 & 764 µs & 2.77 & 121 KB \\
  \tablerule
  \code{Read\allowbreak{}And\allowbreak{}Verify} & 100 & 741 µs & 2.68 & 136 KB \\
  \tablerule
  \code{Read} & 1000 & 846 µs & 1.00 & 269 KB \\
  \tablerule
  \code{Read\allowbreak{}Intercepted} & 1000 & 3858 µs & 4.70 & 1120 KB \\
  \tablerule
  \code{Read\allowbreak{}And\allowbreak{}Verify} & 1000 & 4288 µs & 5.22 & 1268 KB \\
\end{datatable}

The result is surprising: most of the cost is not the check, but having a materialization
interceptor registered at all. EF then produces, for every row, the data that
\code{Materialization\allowbreak{}Interception\allowbreak{}Data} offers, and leaves the fast,
compiled materialization path: at 1000 rows about 0.85 KB and 3 µs per row. The plugin's check
(provider from the cache, read the value, compare) adds around 0.15 KB and 0.4 µs per row. At
ten rows both disappear in the round trip; at a thousand rows reading on the client becomes
about five times as expensive.

\code{Query\allowbreak{}Compilation\allowbreak{}Benchmark} needs no database:
\code{To\allowbreak{}Query\allowbreak{}String} runs the pipeline up to the SQL without a
connection. \code{Compile} empties EF's query cache (the \code{IMemory\allowbreak{}Cache} of
the internal service provider) before every call; \code{Cache\allowbreak{}Hit} runs the same
query shape with a new parameter value. Compared are EF Core alone, the plugin, and the plugin
with its query filter:

\begin{datatable}{H{1.40}H{0.90}H{0.90}H{0.90}H{0.90}}
  \textbf{Context} & \textbf{Compile} & \textbf{Allocated} & \textbf{Cache hit} &
  \textbf{Allocated} \\
  \tablehead
  EF Core & 2.64 ms & 146 KB & 18 µs & 9.9 KB \\
  \tablerule
  Plugin & 2.68 ms & 147 KB & 15 µs & 9.9 KB \\
  \tablerule
  Plugin with query filter & 2.78 ms & 160 KB & 15 µs & 10.2 KB \\
\end{datatable}

Translator plugins and the evaluatable filter cost nothing measurable when compiling; the query filter makes the expression tree, and thus the
translation, a few percent longer. On a cache hit every difference disappears in the noise.
All values come from a short job (.NET 10, author's machine); its error bars are wide, but the
orders of magnitude are clear:

\begin{datatable}{H{0.80}H{0.90}H{1.30}}
  \textbf{Phase} & \textbf{Frequency} & \textbf{Plugin overhead} \\
  \tablehead
  Compilation & Once per query shape & Hardly measurable, a few percent with the query filter \\
  \tablerule
  Cache hit & Every query & Not measurable \\
  \tablerule
  SaveChanges & Every save & About 0.5 µs per entity (40 to 50 percent of change tracking) \\
  \tablerule
  Materialization & Every row read & About 3.5 µs per row with
  \code{Verify\allowbreak{}Materialized\allowbreak{}Tenants(\allowbreak{})}, none without \\
\end{datatable}

A dry run checks that all benchmarks run, without meaningful measurements:

\begin{sourcecode}[All benchmarks, dry run]{text}
dotnet run -c Release --project benchmarks/Acme.EntityFrameworkCore.Tenancy.Benchmarks \
    -- --filter '*' --job dry
\end{sourcecode}

\subsection*{When the check on reading pays off}

\code{Verify\allowbreak{}Materialized\allowbreak{}Tenants(\allowbreak{})} is defense in depth:
it does not protect against an attack that RLS does not stop, but notices when RLS does not
apply at all, for example because the application connects as the owner, a superuser or a role
with \code{BYPASSRLS}, or works on a self-opened connection without
\code{Apply\allowbreak{}Tenant\allowbreak{}Session}. That leads to the recommendation:

\begin{itemize}
  \item \textbf{Turn it on} when a misconfiguration costs more than the CPU time: data with
    high protection requirements, many environments with their own configuration, small
    result sets as typical API queries return them. There the cost disappears in the round
    trip.
  \item \textbf{Turn it off} for read-heavy paths with large result sets, such as exports and
    reports. There the database carries the isolation alone, backed by the role separation
    (page~\pageref{sec:deployment}) and by tests that read with the application's restricted
    role.
  \item \textbf{In between}: always on in tests and staging, in production only for the
    context of the sensitive data. The option applies per context configuration, not
    globally.
\end{itemize}

\begin{tip}
  For the startup of large models, the compiled model
  (chapter~\ref{ch:compiled-models-trimming-and-nativeaot}) is the most effective lever. A
  plugin's conventions then no longer run at run time at all; that is why it pays to test
  compiled model support rather than merely claim it.
\end{tip}

\section{Second-level caches}

A second-level cache (for example \code{EFCore\allowbreak{}Second\allowbreak{}Level\allowbreak{}Cache\allowbreak{}Interceptor},
or results kept in Redis or \code{IMemory\allowbreak{}Cache}) is the obvious next step when
queries are too slow. With multi-tenancy it is also the shortest way to a data leak. The cache
identifies a result by the SQL text or the normalized LINQ expression, and with row level
security the tenant appears in neither: the policy reads it from the session.

\begin{datatable}{H{0.85}H{1.15}}
  \textbf{Step} & \textbf{What happens} \\
  \tablehead
  Tenant A queries & The database filters by the session value, the cache stores the rows
  of A under the SQL text \\
  \tablerule
  Tenant B runs the same query & Same SQL text, same key: the cache returns the rows of A.
  The database is never asked, so the policies never run \\
\end{datatable}

The plugin's other safeguards do not help either. The SaveChanges interceptor only sees
saving. The materialization check (chapter~\ref{ch:materialization-from-rows-to-objects})
only runs when EF builds entities from rows; a cache that stores the entities or DTOs hands
them out without that step. A filter with the tenant as a query parameter (Cosmos DB, test
mode) puts the tenant into the key only if the cache includes parameter values, which not
every cache does.

\begin{important}[The tenant belongs in every cache key]
  A cache in front of a tenant-owned query is only safe if its key contains the current
  tenant, set by the application, not derived from the SQL. Without a tenant nothing is cached.
  Most cache libraries have a hook for the key or a key prefix; where they have none, the cache
  does not belong in front of tenant data. A test with two tenants and the same query shows
  whether the hook works.
\end{important}

\begin{sourcecode}[Example]{csharp}
// Wherever query results are cached: the tenant belongs in every key
public sealed class TenantCacheKeys
{
    private readonly ITenantProvider _tenants;

    public TenantCacheKeys(ITenantProvider tenants) => _tenants = tenants;

    public string For(string queryKey)
        => _tenants.CurrentTenantId is { } tenant
            ? $"tenant:{tenant}:{queryKey}"
            : throw new InvalidOperationException("No tenant: do not cache the result.");
}
\end{sourcecode}

% ------------------------------------------------------------
\chapter{Using the plugin in ASP.NET Core}\label{ch:using-the-plugin-in-asp}
\versionrisk{stable}

The sample \code{samples/\allowbreak{}Shop\allowbreak{}Api} shows the way from the plugin
into a web API. The tenant comes from the claim \code{tenant\_\allowbreak{}id} of the
signed-in user:

% source: samples/ShopApi/HttpTenantProvider.cs
\begin{sourcecode}[samples/\allowbreak{}Shop\allowbreak{}Api/\allowbreak{}Http\allowbreak{}Tenant\allowbreak{}Provider.\allowbreak{}cs]{csharp}
public sealed class HttpTenantProvider : ITenantProvider
{
    public const string ClaimType = "tenant_id";

    private readonly IHttpContextAccessor _httpContextAccessor;

    public HttpTenantProvider(IHttpContextAccessor httpContextAccessor)
    {
        ArgumentNullException.ThrowIfNull(httpContextAccessor);

        _httpContextAccessor = httpContextAccessor;
    }

    // Determined anew on every access: also works as a singleton (pooling)
    public TenantId? CurrentTenantId
    {
        get
        {
            var value = _httpContextAccessor.HttpContext?.User.FindFirst(ClaimType)?.Value;
            return Guid.TryParse(value, out var id) ? new TenantId(id) : null;
        }
    }
}
\end{sourcecode}

% source: samples/ShopApi/Program.cs
\begin{sourcecode}[samples/\allowbreak{}Shop\allowbreak{}Api/\allowbreak{}Program.\allowbreak{}cs (\allowbreak{}excerpt)]{csharp}
builder.Services.AddHttpContextAccessor();
builder.Services.AddAcmeTenancy<HttpTenantProvider>();
builder.Services.AddDbContext<ShopApiContext>((sp, options) => options
    .UseNpgsql(builder.Configuration.GetConnectionString("Shop"))
    .UseAcmeTenancy(sp));

// No tenant filter in code: RLS in the database provides the separation
app.MapGet("/orders", async (ShopApiContext db) =>
    await db.Orders.OrderBy(o => o.Id).ToListAsync());
\end{sourcecode}

\begin{important}[Where the tenant comes from]
  The tenant must come from a trustworthy source: a claim from a signed token. A header or
  query parameter that the client sets freely undermines the entire isolation. The sample
  contains a header simulation for the development environment only and labels it
  explicitly.
\end{important}

For \code{Add\allowbreak{}Db\allowbreak{}Context\allowbreak{}Pool}, register the provider as
a singleton. This is possible because \code{Http\allowbreak{}Tenant\allowbreak{}Provider}
determines the tenant anew on every access; the interceptors ask it on every connection and
every save:

\begin{sourcecode}[Variant with pooling]{csharp}
builder.Services.AddAcmeTenancy<HttpTenantProvider>(ServiceLifetime.Singleton);
builder.Services.AddDbContextPool<ShopApiContext>((sp, options) => options
    .UseNpgsql(connectionString)
    .UseAcmeTenancy(sp));   // sp: root provider; interceptors are created once
\end{sourcecode}

\section{DbContext factories}

\code{Add\allowbreak{}Db\allowbreak{}Context\allowbreak{}Factory} is the usual choice in
Blazor Server and for parallel work. Unlike \code{Add\allowbreak{}Db\allowbreak{}Context},
it registers the factory and the options as \textbf{singletons} by default. The options
callback then receives the root provider, not the scope of the current request. For a scoped
tenant provider that holds state per scope, this has consequences:

\begin{datatable}{H{0.87}H{1.13}}
  \textbf{Registration} & \textbf{Behavior} \\
  \tablehead
  Singleton factory, scope validation on (ASP.NET Core in Development) & Exception: “Cannot
  resolve scoped service 'ITenantProvider' from root provider.” \\
  \tablerule
  Singleton factory, scope validation off (production default) & No error. The interceptors
  get an instance from the root container that no request ever fills: no tenant, saving
  throws, queries see nothing \\
  \tablerule
  Factory registered as \code{Scoped} & Every scope gets its own factory with its own tenant
  provider: correct \\
\end{datatable}

\begin{sourcecode}[Factory with the tenant of the request]{csharp}
builder.Services.AddDbContextFactory<ShopApiContext>(
    (sp, options) => options.UseNpgsql(connectionString).UseAcmeTenancy(sp),
    ServiceLifetime.Scoped);         // factory and options per scope
\end{sourcecode}

A provider that determines the tenant anew on every access, such as
\code{Http\allowbreak{}Tenant\allowbreak{}Provider}, works in every variant; the problem
affects providers with state per scope.
\code{Add\allowbreak{}Pooled\allowbreak{}Db\allowbreak{}Context\allowbreak{}Factory} always
behaves like pooling: the options exist once, so the tenant provider must be such a
singleton. \code{Db\allowbreak{}Context\allowbreak{}Factory\allowbreak{}Tests} shows all
three variants.

\section{Outside of HTTP requests}

Background services, message handlers and scheduled jobs have no request and no claim. The
tenant then comes from the job itself, for example from the message or from a list of
tenants to process. The sample's \code{Ambient\allowbreak{}Tenant\allowbreak{}Provider}
stores it in an \code{Async\allowbreak{}Local}, which flows with the asynchronous control
flow:

% source: samples/ShopSample/AmbientTenantProvider.cs
\begin{sourcecode}[samples/\allowbreak{}Shop\allowbreak{}Sample/\allowbreak{}Ambient\allowbreak{}Tenant\allowbreak{}Provider.\allowbreak{}cs (\allowbreak{}shortened)]{csharp}
public sealed class AmbientTenantProvider : ITenantProvider
{
    private static readonly AsyncLocal<TenantId?> Current = new();

    public TenantId? CurrentTenantId => Current.Value;

    public static IDisposable Use(TenantId tenant)
    {
        var previous = Current.Value;
        Current.Value = tenant;
        return new Restore(previous);     // sets the previous tenant back on Dispose
    }

    // private sealed class Restore : IDisposable { ... Current.Value = _previous; }
}
\end{sourcecode}

\begin{sourcecode}[Job pattern (\allowbreak{}tests/\allowbreak{}.\allowbreak{}.\allowbreak{}.\allowbreak{}/\allowbreak{}Background\allowbreak{}Tenant\allowbreak{}Tests.\allowbreak{}cs,\allowbreak{} simplified)]{csharp}
// Tenant first, then scope and context: the interceptors are created inside the scope
using (AmbientTenantProvider.Use(tenant))
{
    await using var scope = scopeFactory.CreateAsyncScope();
    var db = scope.ServiceProvider.GetRequiredService<ShopContext>();
    // ... work for this tenant ...
}
\end{sourcecode}

Because the provider reads the \code{Async\allowbreak{}Local} on every access, it can be
registered as a singleton. Parallel jobs keep their own tenant, and after the \code{using}
block no tenant is left behind; \code{Background\allowbreak{}Tenant\allowbreak{}Tests}
checks both with 20 parallel jobs. An application with web requests \emph{and} jobs combines
both sources in one provider: the claim first, otherwise the ambient tenant.

\textbf{Messages.} With a message broker (RabbitMQ, Kafka, Azure Service Bus) or a framework on
top of it (MassTransit, Wolverine), the sender writes the current tenant into a header of the
message, and the receiving side sets it before the consumer runs. The place for that is the
framework's consume filter or middleware, so that no single consumer can forget it:

\begin{sourcecode}[Consume filter (\allowbreak{}simplified,\allowbreak{} independent of the framework)]{csharp}
// Wherever the broker hands over a message: tenant from the header, then scope and context
public async Task HandleAsync(BrokerMessage message, CancellationToken cancellationToken)
{
    if (!message.Headers.TryGetValue("X-Tenant-Id", out var value)
        || !Guid.TryParse(value, out var id))
        throw new InvalidOperationException("Message without a tenant: rejected.");

    using (AmbientTenantProvider.Use(new TenantId(id)))
    {
        await using var scope = scopeFactory.CreateAsyncScope();
        var handler = scope.ServiceProvider.GetRequiredService<IOrderHandler>();
        await handler.HandleAsync(message, cancellationToken);   // uses ShopContext
    }
}
\end{sourcecode}

Two rules carry over from the jobs. The header is read before the scope is created, because
the context and its interceptors live in the scope. And a message without a valid tenant is
rejected and ends up in the dead letter queue; it is never processed for all tenants. A scoped
tenant provider with a setter works the same way, as long as the filter sets the tenant in the
consumer's scope before the context is resolved.

\begin{important}[Never fall back to “all tenants”]
  Without a tenant, a job must fail or see nothing, never everything. The plugin behaves
  this way: the session value stays empty, the policies match no rows and saving throws.
  Jobs that really need all tenants use the reporting role from
  chapter~\ref{ch:row-level-security-in-production}.
\end{important}

% ------------------------------------------------------------
\chapter{CI and API stability}\label{ch:ci-and-api-stability}
\versionrisk{stable}

\section{The workflow}

The GitHub Actions workflow in the repository implements the recommendations from
chapters~\ref{ch:test-strategy} and \ref{ch:versioning-and-compatibility}. The job
\code{build} builds and tests without and with a database, \code{handbook} builds the
handbook, \code{consumer} checks the packages as a consumer. The job \code{cosmos} (emulator,
not blocking yet) is left out of the excerpt:

\begin{sourcecode}[.\allowbreak{}github/\allowbreak{}workflows/\allowbreak{}ci.\allowbreak{}yml (\allowbreak{}shortened)]{text}
jobs:
  build:
    runs-on: ubuntu-latest
    steps:
      - uses: actions/checkout@v4
      - uses: actions/setup-dotnet@v4
        with: { global-json-file: global.json }
      - run: dotnet build -c Release
      - run: dotnet test tests/Acme.EntityFrameworkCore.Tenancy.Tests -c Release --no-build
      - run: dotnet test tests/Acme.EntityFrameworkCore.Tenancy.IntegrationTests
          -c Release --no-build
      - run: dotnet test tests/Acme.EntityFrameworkCore.Tenancy.SqlServer.Tests
          -c Release --no-build

  handbook:                            # both PDFs with LuaLaTeX, in docs/handbook
    runs-on: ubuntu-latest
    container: texlive/texlive@sha256:540d…   # TeX Live 2026, pinned by digest
    steps:
      # ... checkout ...
      - run: python3 check_code.py       # line length, excerpts against the sources
      - run: python3 check_structure.py  # editions in sync, references, tables
      - run: latexmk EF-Core-Plugin-Handbuch.tex
      - run: latexmk EF-Core-Plugin-Handbook.tex

  consumer:                            # check the packages as a consumer
    steps:
      - run: eng/consumer-test/run.sh --nativeaot
\end{sourcecode}

Code excerpts go stale silently when the code changes. A comment line
\code{\% source: <path>} directly before the listing therefore binds an excerpt to its source
file, and \code{check\_\allowbreak{}code.\allowbreak{}py} compares the two: comments and elisions
(\code{...}) do not count, so the German edition may translate comments and excerpts may be
shortened; every other line must appear in the source in the same order, line breaks may
differ. With the suffix \code{(strings)} translated messages are left out as well. An excerpt
whose title names a file of the repository and that is not bound must be called “simplified”.
\code{check\_\allowbreak{}structure.\allowbreak{}py} checks that both editions have the same
sections, listings, bindings, tables and callouts per chapter, that every reference target exists
and that the column widths of every table add up.

The sample builds against stable versions only (.NET 10, EF Core 10). A second workflow,
\code{compatibility.\allowbreak{}yml}, runs weekly and tests against the newest stable 10.0.x
patches, because the regular CI builds against the beginning of the version range. The
channels and a sketch for the next EF major version with previews and RCs are described in
section “Patches and previews in CI” of chapter~\ref{ch:versioning-and-compatibility}.

\section{Protecting the public API}

A plugin has two public APIs: the one for users (Fluent API, options, types) and the one for
extenders (virtual methods, dependencies, public services). Both break easily and unnoticed.
The
\code{Microsoft.\allowbreak{}Code\allowbreak{}Analysis.\allowbreak{}Public\allowbreak{}Api\allowbreak{}Analyzers}
package keeps the public API in two text files
(\code{Public\allowbreak{}API.\allowbreak{}Shipped.\allowbreak{}txt},
\code{Public\allowbreak{}API.\allowbreak{}Unshipped.\allowbreak{}txt}) and reports every
deviation. Every API change thus becomes visible as a diff in the pull request.

\begin{note}[Setup]
  On the first build, the analyzer reports every public API as RS0016 because the files are
  empty. The code fix “Add all items … to the public API” fills them, on the command line with
  \code{dotnet format analyzers -{}-diagnostics RS0016 -{}-severity warn}. At release time,
  the content moves from Unshipped to Shipped; from then on, removing a shipped API is a
  visible breaking change.
\end{note}

\begin{important}[Warnings are not enough]
  As long as RS0016 and RS0017 are only warnings, they get lost in the build output. In the
  sample, the files thus stayed empty up to edition~1.0: 178 warnings in a full build, none
  in incremental ones, and nobody noticed.
  \code{Directory.\allowbreak{}Build.\allowbreak{}props} therefore turns both into errors
  (\code{Warnings\allowbreak{}As\allowbreak{}Errors}); a new public API now has to be entered
  in \code{Public\allowbreak{}API.\allowbreak{}Unshipped.\allowbreak{}txt}, otherwise the
  build breaks.
\end{important}
